%=====================================
%   20250802.tex
%   20250802 
%   toshi2019
%=====================================
% ------------------------
% documentclass
% ------------------------
\documentclass[leqno, 10pt, a4paper, dvipdfmx]{jsarticle}
\title{ゼミノート（08月02日分）}
\author{Toshi2019}
\date{\today}

% ------------------------
% usepackage
% ------------------------
\usepackage{algorithm}
\usepackage{algorithmic}
\usepackage{amscd}
\usepackage{amsfonts}
\usepackage{amsmath}
\usepackage[psamsfonts]{amssymb}
\usepackage{amsthm}
\usepackage{ascmac}
\usepackage{bm}
\usepackage{color}
\usepackage{enumerate}
\usepackage{fancybox}
\usepackage[stable]{footmisc}
\usepackage{graphicx}
\usepackage{listings}
\usepackage{mathrsfs}
\usepackage{mathtools}
\usepackage{otf}
\usepackage{pifont}
\usepackage{proof}
\usepackage{subfigure}
\usepackage{tikz}
\usepackage{verbatim}
\usepackage[all]{xy}

\usetikzlibrary{cd}
\usetikzlibrary{arrows.meta}
\usetikzlibrary{calc}

% ================================
% パッケージを追加する場合のスペース 
\usepackage[dvipdfmx]{hyperref}
\usepackage{xcolor}
\definecolor{darkgreen}{rgb}{0,0.45,0} 
\definecolor{darkred}{rgb}{0.75,0,0}
\definecolor{darkblue}{rgb}{0,0,0.6} 
\hypersetup{
    colorlinks=true,
    citecolor=darkgreen,
    linkcolor=darkred,
    urlcolor=darkblue,
}
\usepackage{pxjahyper}

%\usepackage[mathscr]{euscript} % mathscr のフォントを変更

%\usepackage{mmasym}

%=================================


% --------------------------
% theoremstyle
% --------------------------
\theoremstyle{definition}

% --------------------------
% newtheoem
% --------------------------

% 日本語で定理, 命題, 証明などを番号付きで用いるためのコマンドです. 
% If you want to use theorem environment in Japanece, 
% you can use these code. 
% Attention!
% All theorem enivironment numbers depend on 
% only section numbers.
\newtheorem{Axiom}{公理}[section]
\newtheorem{Definition}[Axiom]{定義}
\newtheorem{Theorem}[Axiom]{定理}
\newtheorem{Proposition}[Axiom]{命題}
\newtheorem{Lemma}[Axiom]{補題}
\newtheorem{Corollary}[Axiom]{系}
\newtheorem{Example}[Axiom]{例}
\newtheorem{Claim}[Axiom]{主張}
\newtheorem{Property}[Axiom]{性質}
\newtheorem{Attention}[Axiom]{注意}
\newtheorem{Question}[Axiom]{問}
\newtheorem{Problem}[Axiom]{問題}
\newtheorem{Consideration}[Axiom]{考察}
\newtheorem{Alert}[Axiom]{警告}
\newtheorem{Fact}[Axiom]{事実}


% 日本語で定理, 命題, 証明などを番号なしで用いるためのコマンドです. 
% If you want to use theorem environment with no number in Japanese, You can use these code.
\newtheorem*{Axiom*}{公理}
\newtheorem*{Definition*}{定義}
\newtheorem*{Theorem*}{定理}
\newtheorem*{Proposition*}{命題}
\newtheorem*{Lemma*}{補題}
\newtheorem*{Example*}{例}
\newtheorem*{Corollary*}{系}
\newtheorem*{Claim*}{主張}
\newtheorem*{Property*}{性質}
\newtheorem*{Attention*}{注意}
\newtheorem*{Question*}{問}
\newtheorem*{Problem*}{問題}
\newtheorem*{Consideration*}{考察}
\newtheorem*{Alert*}{警告}
\newtheorem{Fact*}{事実}


% 英語で定理, 命題, 証明などを番号付きで用いるためのコマンドです. 
%　If you want to use theorem environment in English, You can use these code.
%all theorem enivironment number depend on only section number.
\newtheorem{Axiom+}{Axiom}[section]
\newtheorem{Definition+}[Axiom+]{Definition}
\newtheorem{Theorem+}[Axiom+]{Theorem}
\newtheorem{Proposition+}[Axiom+]{Proposition}
\newtheorem{Lemma+}[Axiom+]{Lemma}
\newtheorem{Example+}[Axiom+]{Example}
\newtheorem{Corollary+}[Axiom+]{Corollary}
\newtheorem{Claim+}[Axiom+]{Claim}
\newtheorem{Property+}[Axiom+]{Property}
\newtheorem{Attention+}[Axiom+]{Attention}
\newtheorem{Question+}[Axiom+]{Question}
\newtheorem{Problem+}[Axiom+]{Problem}
\newtheorem{Consideration+}[Axiom+]{Consideration}
\newtheorem{Alert+}{Alert}
\newtheorem{Fact+}[Axiom+]{Fact}
\newtheorem{Remark+}[Axiom+]{Remark}

% ----------------------------
% commmand
% ----------------------------
% 執筆に便利なコマンド集です. 
% コマンドを追加する場合は下のスペースへ. 

% 集合の記号 (黒板文字)
\newcommand{\NN}{\mathbb{N}}
\newcommand{\ZZ}{\mathbb{Z}}
\newcommand{\QQ}{\mathbb{Q}}
\newcommand{\RR}{\mathbb{R}}
\newcommand{\CC}{\mathbb{C}}
\newcommand{\PP}{\mathbb{P}}
\newcommand{\KK}{\mathbb{K}}


% 集合の記号 (太文字)
\newcommand{\nn}{\mathbf{N}}
\newcommand{\zz}{\mathbf{Z}}
\newcommand{\qq}{\mathbf{Q}}
\newcommand{\rr}{\mathbf{R}}
\newcommand{\cc}{\mathbf{C}}
\newcommand{\pp}{\mathbf{P}}
\newcommand{\kk}{\mathbf{k}}

% 特殊な写像の記号
\newcommand{\ev}{\mathop{\mathrm{ev}}\nolimits} % 値写像
\newcommand{\pr}{\mathop{\mathrm{pr}}\nolimits} % 射影

% スクリプト体にするコマンド
%   例えば　{\mcal C} のように用いる
\newcommand{\mcal}{\mathcal}

% 花文字にするコマンド 
%   例えば　{\h C} のように用いる
\newcommand{\h}{\mathscr}

% ヒルベルト空間などの記号
\newcommand{\F}{\mcal{F}}
\newcommand{\X}{\mcal{X}}
\newcommand{\Y}{\mcal{Y}}
\newcommand{\Hil}{\mcal{H}}
\newcommand{\RKHS}{\Hil_{k}}
\newcommand{\Loss}{\mcal{L}_{D}}
\newcommand{\MLsp}{(\X, \Y, D, \Hil, \Loss)}

% 偏微分作用素の記号
\newcommand{\p}{\partial}

% 角カッコの記号 (内積は下にマクロがあります)
\newcommand{\lan}{\langle}
\newcommand{\ran}{\rangle}



% 圏の記号など
\newcommand{\Set}{{\bf Set}}
\newcommand{\Vect}{{\bf Vect}}
\newcommand{\FDVect}{{\bf FDVect}}
%\newcommand{\Ring}{{\bf Ring}}
\newcommand{\Ab}{{\bf Ab}}
\newcommand{\Mod}{\mathop{\mathrm{Mod}}\nolimits}
\newcommand{\Modf}{\mathop{\mathrm{Mod}^\mathrm{f}}\nolimits}
\newcommand{\CGA}{{\bf CGA}}
\newcommand{\GVect}{{\bf GVect}}
\newcommand{\Lie}{{\bf Lie}}
\newcommand{\dLie}{{\bf Liec}}



% 射の集合など
\newcommand{\Map}{\mathop{\mathrm{Map}}\nolimits} % 写像の集合
\newcommand{\Hom}{\mathop{\mathrm{Hom}}\nolimits} % 射集合
\newcommand{\End}{\mathop{\mathrm{End}}\nolimits} % 自己準同型の集合
\newcommand{\Aut}{\mathop{\mathrm{Aut}}\nolimits} % 自己同型の集合
\newcommand{\Mor}{\mathop{\mathrm{Mor}}\nolimits} % 射集合
\newcommand{\Ker}{\mathop{\mathrm{Ker}}\nolimits} % 核
\newcommand{\Img}{\mathop{\mathrm{Im}}\nolimits} % 像
\newcommand{\Cok}{\mathop{\mathrm{Coker}}\nolimits} % 余核
\newcommand{\Cim}{\mathop{\mathrm{Coim}}\nolimits} % 余像

% その他便利なコマンド
\newcommand{\dip}{\displaystyle} % 本文中で数式モード
\newcommand{\e}{\varepsilon} % イプシロン
\newcommand{\dl}{\delta} % デルタ
\newcommand{\pphi}{\varphi} % ファイ
\newcommand{\ti}{\tilde} % チルダ
\newcommand{\pal}{\parallel} % 平行
\newcommand{\op}{{\rm op}} % 双対を取る記号
\newcommand{\lcm}{\mathop{\mathrm{lcm}}\nolimits} % 最小公倍数の記号
\newcommand{\Probsp}{(\Omega, \F, \P)} 
\newcommand{\argmax}{\mathop{\rm arg~max}\limits}
\newcommand{\argmin}{\mathop{\rm arg~min}\limits}





% ================================
% コマンドを追加する場合のスペース 
%\newcommand{\OO}{\mcal{O}}



\makeatletter
\renewenvironment{proof}[1][\proofname]{\par
  \pushQED{\qed}%
  \normalfont \topsep6\p@\@plus6\p@\relax
  \trivlist
  \item[\hskip\labelsep
%        \itshape
         \bfseries
%    #1\@addpunct{.}]\ignorespaces
    {#1}]\ignorespaces
}{%
  \popQED\endtrivlist\@endpefalse
}
\makeatother

\renewcommand{\proofname}{証明.}



%\renewcommand\proofname{\bf 証明} % 証明
\numberwithin{equation}{section}
\newcommand{\cTop}{\textsf{Top}}
%\newcommand{\cOpen}{\textsf{Open}}
\newcommand{\Op}{\mathop{\textsf{Op}}\nolimits}
\newcommand{\Ob}{\mathop{\textrm{Ob}}\nolimits}
\newcommand{\id}{\mathop{\mathrm{id}}\nolimits}
\newcommand{\pt}{\mathop{\mathrm{pt}}\nolimits}
\newcommand{\res}{\mathop{\rho}\nolimits}
\newcommand{\A}{\mcal{A}}
\newcommand{\B}{\mcal{B}}
\newcommand{\C}{\mcal{C}}
\newcommand{\D}{\mcal{D}}
\newcommand{\E}{\mcal{E}}
\newcommand{\G}{\mcal{G}}
%\newcommand{\H}{\mcal{H}}
\newcommand{\I}{\mcal{I}}
\newcommand{\J}{\mcal{J}}
\newcommand{\OO}{\mcal{O}}
\newcommand{\Ring}{\mathop{\textsf{Ring}}\nolimits}
\newcommand{\cAb}{\mathop{\textsf{Ab}}\nolimits}
%\newcommand{\Ker}{\mathop{\mathrm{Ker}}\nolimits}
\newcommand{\im}{\mathop{\mathrm{Im}}\nolimits}
\newcommand{\Coker}{\mathop{\mathrm{Coker}}\nolimits}
\newcommand{\Coim}{\mathop{\mathrm{Coim}}\nolimits}
\newcommand{\rank}{\mathop{\mathrm{rank}}\nolimits}
\newcommand{\Ht}{\mathop{\mathrm{Ht}}\nolimits}
\newcommand{\supp}{\mathop{\mathrm{supp}}\nolimits}
\newcommand{\colim}{\mathop{\mathrm{colim}}}
\newcommand{\Tor}{\mathop{\mathrm{Tor}}\nolimits}

\newcommand{\cat}{\mathscr{C}}

%筆記体
\newcommand{\cA}{\mcal{A}}
\newcommand{\cB}{\mcal{B}}
\newcommand{\cC}{\mcal{C}}
\newcommand{\cD}{\mcal{D}}
\newcommand{\cE}{\mcal{E}}
\newcommand{\cF}{\mcal{F}}
\newcommand{\cG}{\mcal{G}}
\newcommand{\cH}{\mcal{H}}
\newcommand{\cI}{\mcal{I}}
\newcommand{\cJ}{\mcal{J}}
\newcommand{\cK}{\mcal{K}}
\newcommand{\cL}{\mcal{L}}
\newcommand{\cM}{\mcal{M}}
\newcommand{\cN}{\mcal{N}}
\newcommand{\cO}{\mcal{O}}
\newcommand{\cP}{\mcal{P}}
\newcommand{\cQ}{\mcal{Q}}
\newcommand{\cR}{\mcal{R}}
\newcommand{\cS}{\mcal{S}}
\newcommand{\cT}{\mcal{T}}
\newcommand{\cU}{\mcal{U}}
\newcommand{\cV}{\mcal{V}}
\newcommand{\cW}{\mcal{W}}
\newcommand{\cX}{\mcal{X}}
\newcommand{\cY}{\mcal{Y}}
\newcommand{\cZ}{\mcal{Z}}


\newcommand{\scA}{\mathscr{A}}
\newcommand{\scB}{\mathscr{B}}
\newcommand{\scC}{\mathscr{C}}
\newcommand{\scD}{\mathscr{D}}
\newcommand{\scE}{\mathscr{E}}
\newcommand{\scF}{\mathscr{F}}
\newcommand{\scN}{\mathscr{N}}
\newcommand{\scO}{\mathscr{O}}
\newcommand{\scV}{\mathscr{V}}
\newcommand{\scU}{\mathscr{U}}


\newcommand{\ibA}{\mathop{\text{\textit{\textbf{A}}}}}
\newcommand{\ibB}{\mathop{\text{\textit{\textbf{B}}}}}
\newcommand{\ibC}{\mathop{\text{\textit{\textbf{C}}}}}
\newcommand{\ibD}{\mathop{\text{\textit{\textbf{D}}}}}
\newcommand{\ibE}{\mathop{\text{\textit{\textbf{E}}}}}
\newcommand{\ibF}{\mathop{\text{\textit{\textbf{F}}}}}
\newcommand{\ibG}{\mathop{\text{\textit{\textbf{G}}}}}
\newcommand{\ibH}{\mathop{\text{\textit{\textbf{H}}}}}
\newcommand{\ibI}{\mathop{\text{\textit{\textbf{I}}}}}
\newcommand{\ibJ}{\mathop{\text{\textit{\textbf{J}}}}}
\newcommand{\ibK}{\mathop{\text{\textit{\textbf{K}}}}}
\newcommand{\ibL}{\mathop{\text{\textit{\textbf{L}}}}}
\newcommand{\ibM}{\mathop{\text{\textit{\textbf{M}}}}}
\newcommand{\ibN}{\mathop{\text{\textit{\textbf{N}}}}}
\newcommand{\ibO}{\mathop{\text{\textit{\textbf{O}}}}}
\newcommand{\ibP}{\mathop{\text{\textit{\textbf{P}}}}}
\newcommand{\ibQ}{\mathop{\text{\textit{\textbf{Q}}}}}
\newcommand{\ibR}{\mathop{\text{\textit{\textbf{R}}}}}
\newcommand{\ibS}{\mathop{\text{\textit{\textbf{S}}}}}
\newcommand{\ibT}{\mathop{\text{\textit{\textbf{T}}}}}
\newcommand{\ibU}{\mathop{\text{\textit{\textbf{U}}}}}
\newcommand{\ibV}{\mathop{\text{\textit{\textbf{V}}}}}
\newcommand{\ibW}{\mathop{\text{\textit{\textbf{W}}}}}
\newcommand{\ibX}{\mathop{\text{\textit{\textbf{X}}}}}
\newcommand{\ibY}{\mathop{\text{\textit{\textbf{Y}}}}}
\newcommand{\ibZ}{\mathop{\text{\textit{\textbf{Z}}}}}

\newcommand{\ibx}{\mathop{\text{\textit{\textbf{x}}}}}

%\newcommand{\Comp}{\mathop{\mathrm{C}}\nolimits}
%\newcommand{\Komp}{\mathop{\mathrm{K}}\nolimits}
%\newcommand{\Domp}{\mathop{\mathsf{D}}\nolimits}%複体のホモトピー圏
%\newcommand{\Comp}{\mathrm{C}}
%\newcommand{\Komp}{\mathrm{K}}
%\newcommand{\Domp}{\mathsf{D}}%複体のホモトピー圏

\newcommand{\Comp}{\mathop{\mathrm{C}}\nolimits}
\newcommand{\Komp}{\mathop{\mathsf{K}}\nolimits}
\newcommand{\Domp}{\mathord{\mathsf{D}}}%\nolimits}
\newcommand{\Kompl}{\mathop{\mathsf{K}^\mathrm{+}}\nolimits}
\newcommand{\Kompu}{\mathop{\mathsf{K}^\mathrm{-}}\nolimits}
\newcommand{\Kompb}{\mathop{\mathsf{K}^\mathrm{b}}\nolimits}
\newcommand{\Dompl}{\mathop{\mathsf{D}^\mathrm{+}}\nolimits}
\newcommand{\Dompu}{\mathop{\mathsf{D}^\mathrm{-}}\nolimits}
\newcommand{\Dompb}{\mathop{\mathsf{D}^\mathrm{b}}\nolimits}
\newcommand{\Dbconic}{\mathop{\mathsf{D}^\mathrm{b}_{\rrp}}\nolimits}
\newcommand{\Dompbf}{\mathop{\mathsf{D}_\mathrm{f}^\mathrm{b}}\nolimits}
\newcommand{\Domplb}{\mathop{\mathsf{D}^\mathrm{lb}}\nolimits}




\newcommand{\CCat}{\Comp(\cat)}
\newcommand{\KCat}{\Komp(\cat)}
\newcommand{\DCat}{\Domp(\cat)}%圏Cの複体のホモトピー圏
%\newcommand{\HOM}{\mathop{\mathscr{H}\hspace{-2pt}om}\nolimits}%内部Hom
\newcommand{\HOM}{\mathop{\mathcal{H}\hspace{-0pt}om}\nolimits}%内部Hom
\newcommand{\RHOM}{\mathop{\mathrm{R}\hspace{-1.5pt}\HOM}\nolimits}

\newcommand{\muS}{\mathop{\mathrm{SS}}\nolimits}
\newcommand{\RG}{\mathop{\mathrm{R}\hspace{-0pt}\Gamma}\nolimits}
\newcommand{\RHom}{\mathop{\mathrm{R}\hspace{-1.5pt}\Hom}\nolimits}
\newcommand{\Rder}{\mathrm{R}}

\newcommand{\simar}{\mathrel{\overset{\sim}{\rightarrow}}}%同型右矢印
\newcommand{\simarr}{\mathrel{\overset{\sim}{\longrightarrow}}}%同型右矢印
\newcommand{\simra}{\mathrel{\overset{\sim}{\leftarrow}}}%同型左矢印
\newcommand{\simrra}{\mathrel{\overset{\sim}{\longleftarrow}}}%同型左矢印

\newcommand{\hocolim}{{\mathrm{hocolim}}}
\newcommand{\indlim}[1][]{\mathop{\varinjlim}\limits_{#1}}
\newcommand{\sindlim}[1][]{\smash{\mathop{\varinjlim}\limits_{#1}}\,}
\newcommand{\Pro}{\mathrm{Pro}}
\newcommand{\Ind}{\mathrm{Ind}}
\newcommand{\prolim}[1][]{\mathop{\varprojlim}\limits_{#1}}
\newcommand{\sprolim}[1][]{\smash{\mathop{\varprojlim}\limits_{#1}}\,}

\newcommand{\Sh}{\mathrm{Sh}}
\newcommand{\PSh}{\mathrm{PSh}}

\newcommand{\rmD}{\mathsf{D}}
\newcommand{\vD}{\mathbf{D}}

\newcommand{\Lloc}[1][]{\mathord{\mathcal{L}^1_{\mathrm{loc},{#1}}}}
\newcommand{\ori}{\mathord{\mathrm{or}}}
\newcommand{\Db}{\mathord{\mathscr{D}b}}

\newcommand{\codim}{\mathop{\mathrm{codim}}\nolimits}



\newcommand{\gld}{\mathop{\mathrm{gld}}\nolimits}
\newcommand{\wgld}{\mathop{\mathrm{wgld}}\nolimits}


\newcommand{\tens}[1][]{\mathbin{\otimes_{\raise1.5ex\hbox to-.1em{}{#1}}}}
\newcommand{\etens}{\mathbin{\boxtimes}}
\newcommand{\ltens}[1][]{\mathbin{\overset{\mathrm{L}}\tens}_{#1}}
\newcommand{\mtens}[1][]{\mathbin{\overset{\mathrm{\mu}}\tens}_{#1}}
\newcommand{\lltens}[1][]{\mathbin{{\mathop{\tens}\limits^{\mathrm{L}}_{#1}}}}
\newcommand{\lletens}[1][]{\mathbin{{\mathop{\boxtimes}\limits^{\mathrm{L}}_{#1}}}}
\newcommand{\letens}{\overset{\mathrm{L}}{\etens}}
\newcommand{\detens}{\underline{\etens}}
\newcommand{\ldetens}{\overset{\mathrm{L}}{\underline{\etens}}}
\newcommand{\dtens}[1][]{{\overset{\mathrm{L}}{\underline{\otimes}}}_{#1}}

\newcommand{\blk}{\mathord{\ \cdot\ }}
\newcommand{\mres}[2][]{{\left.{#1}\right\rvert}_{#2}}


%\newcommand{\hocolim}{{\mathrm{hocolim}}}
%\newcommand{\indlim}[1][]{\mathop{\varinjlim}\limits_{#1}}
%\newcommand{\sindlim}[1][]{\smash{\mathop{\varinjlim}\limits_{#1}}\,}
%\newcommand{\Pro}{\mathrm{Pro}}
%\newcommand{\Ind}{\mathrm{Ind}}
%\newcommand{\prolim}[1][]{\mathop{\varprojlim}\limits_{#1}}
%\newcommand{\sprolim}[1][]{\smash{\mathop{\varprojlim}\limits_{#1}}\,}
\newcommand{\proolim}[1][]{\mathop{\text{\rm``{$\varprojlim$}''}}\limits_{#1}}
\newcommand{\sproolim}[1][]{\smash{\mathop{\rm``{\varprojlim}''}\limits_{#1}}}
\newcommand{\inddlim}[1][]{\mathop{\text{\rm``{$\varinjlim$}''}}\limits_{#1}}
\newcommand{\sinddlim}[1][]{\smash{\mathop{\text{\rm``{$\varinjlim$}''}}\limits_{#1}}\,}
\newcommand{\ooplus}{\mathop{\text{\rm``{$\oplus$}''}}\limits}
\newcommand{\bbigsqcup}{\mathop{``\bigsqcup''}\limits}
\newcommand{\bsqcup}{\mathop{``\sqcup''}\limits}
\newcommand{\dsum}[1][]{\mathbin{\oplus_{#1}}}

\newcommand{\Fct}{\mathop{\mathsf{Fct}}\nolimits}

\newcommand{\pb}[1][]{\mathbin{\mathop{\times}\limits_{#1}}}
\newcommand{\dpb}[1][]{\mathbin{\mathop{\times}\limits_{#1}}}
\newcommand{\tpb}[1][]{\mathbin{\mathop{\times}\nolimits_{#1}}}





%================================================
% 自前の定理環境
%   https://mathlandscape.com/latex-amsthm/
% を参考にした
\newtheoremstyle{mystyle}%   % スタイル名
    {5pt}%                   % 上部スペース
    {5pt}%                   % 下部スペース
    {}%              % 本文フォント
    {}%                  % 1行目のインデント量
    {\bfseries}%                      % 見出しフォント
    {.}%                     % 見出し後の句読点
    {12pt}%                     % 見出し後のスペース
    {\thmname{#1}\thmnumber{ #2}\thmnote{{\hspace{2pt}\normalfont (#3)}}}% % 見出しの書式

\theoremstyle{mystyle}
\newtheorem{AXM}{公理}[section]
\newtheorem{DFN}[AXM]{定義}
\newtheorem{THM}[AXM]{定理}
\newtheorem*{THM*}{定理}
\newtheorem{PRP}[AXM]{命題}
\newtheorem{LMM}[AXM]{補題}
\newtheorem{CRL}[AXM]{系}
\newtheorem{EG}[AXM]{例}
\newtheorem*{EG*}{例}
\newtheorem{RMK}[AXM]{注意}
\newtheorem{CNV}[AXM]{約束}
\newtheorem{CMT}[AXM]{コメント}
\newtheorem*{CMT*}{コメント}
\newtheorem{NTN}[AXM]{記号}

% 定理環境ここまで
%====================================================

\usepackage{framed}
\definecolor{lightgray}{rgb}{0.75,0.75,0.75}
\renewenvironment{leftbar}{%
  \def\FrameCommand{\textcolor{lightgray}{\vrule width 4pt} \hspace{10pt}}% 
  \MakeFramed {\advance\hsize-\width \FrameRestore}}%
{\endMakeFramed}
\newenvironment{redleftbar}{%
  \def\FrameCommand{\textcolor{lightgray}{\vrule width 1pt} \hspace{10pt}}% 
  \MakeFramed {\advance\hsize-\width \FrameRestore}}%
 {\endMakeFramed}



\renewcommand{\Re}{\mathop{\mathrm{Re}}\nolimits}

% =================================





% ---------------------------
% new definition macro
% ---------------------------
% 便利なマクロ集です

% 内積のマクロ
%   例えば \inner<\pphi | \psi> のように用いる
\def\inner<#1>{\langle #1 \rangle}

% ================================
% マクロを追加する場合のスペース 

%=================================



% 位相空間
\newcommand{\Int}[1][]{\mathop{\mathrm{Int}}\nolimits_{#1}}
\newcommand{\Cl}[1][]{\mathop{\mathrm{Cl}}\nolimits_{#1}}
\newcommand{\Bd}[1][]{\mathop{\mathrm{Bd}}\nolimits_{#1}}
\newcommand{\bd}{\mathord{\partial}}
\newcommand{\Fr}[1][]{\mathop{\mathrm{Fr}}\nolimits_{#1}}

\newcommand{\dom}[1][]{\mathop{\mathrm{dom}}\nolimits_{#1}}


\newcommand{\mtan}[1][]{T\hspace{-1.75pt}{#1}}
\newcommand{\mtp}[1][]{{}^{t\!}{#1}} % transpose of the morphism
%\newcommand{\mpb}[1][]{{}^{t\!}{#1}'} % pull-back of the morphism
\newcommand{\mpb}[1][]{{#1}_{\pi}} % pull-back of the morphism
\newcommand{\mdiff}[1][]{{#1}_{d}} % pull-back of the morphism
\newcommand{\mptan}[1][]{{#1}_{\tau}} % pull-back of the morphism


%\newcommand{\pb}[1][]{\mathbin{\mathop{\times}\limits_{#1}}}

\newcommand{\cotb}[1][]{T^{\ast}{#1}}
\newcommand{\conm}[2][]{T_{#1}^{\hspace{0.7pt}\ast}{#2}}
\newcommand{\norb}[2][]{T_{#1}{#2}}

\newcommand{\cotz}[1][]{T^{\ast}_{#1}{#1}}

\newcommand{\cotn}[1][]{\dot{T}^{\ast}{#1}}


\newcommand{\muhom}{\mu hom}
\newcommand{\muHom}[1][]{\mathrm{Hom}^\mu_{\raise1.5ex\hbox to.1em{}#1}}

\newcommand{\mucat}[2][]{\mathsf{D}^{\mathrm{b}}_{#1}(#2)}
%\newcommand{\mucat}[2][]{\mathsf{D}^{\mathrm{b}}_{#1}(#2)}

\newcommand{\conv}[1][]{\mathbin{\mathop{\circ}\limits_{#1}}}

\newcommand{\torus}{\mathbf{T}}

\newcommand{\diag}[1][]{\Delta_{#1}}
%  \[\mu_{\Delta_{Y}}\mathrm{R}\mathcal{H}om(G,F)\]
\newcommand{\micro}[1]{\mu_{#1}\hspace{-2pt}}

\newcommand{\qis}{\stackrel{\text{qis}}{\sim}}%同型右矢印

\newcommand{\Dlcon}{\mathop{\mathsf{D}^{+}_{\rr_{>0}}}\nolimits}
%\newcommand{\Dlcon}{\mathop{\mathsf{D}^{+}_{\rr^+}}\nolimits}
\newcommand{\Dbcon}{\mathop{\mathsf{D}^{\mathrm{b}}_{\rr^+}}\nolimits}



% ----------------------------
% documenet 
% ----------------------------
% 以下, 本文の執筆スペースです. 
% Your main code must be written between 
% begin document and end document.
% ---------------------------

\begin{document}
\maketitle

\section{フーリエ・佐藤変換}
まず錐状層を定義する．
そのために作用つきの空間を考える．
\(X\)を局所コンパクト空間で\(\rr^+\)の作用\(\mu\)が入っているものとする．
つまり，連続写像\(\mu\colon X\times \rr^+\to X\)で
\begin{align*}
    \mu(x,t_1t_2)&=\mu(\mu(x,t_1),t_2)\\
    \mu(x,1)&=x
\end{align*}
をみたすものが与えられているとする．

\begin{leftbar}
\begin{DFN}[{\cite[Definition 3.7.1]{KS90}}]
    \begin{enumerate}[(i)]
        \item 層\(F\in\Mod(A_X)\)が\textbf{錐状} (conic) である
        とは，
        \(X\)の各軌道\(b\)への制限\(F\rvert_b\)が
        局所定数層であることをいう．
        錐状層からなる\(\Mod(A_X)\)の
        充満部分圏を\(\Mod_{\rr^+}(A_X)\)で表す．
        \item 各\(j\in\zz\)に対して\(H^j(F)\)が
        錐状層である対象からなる\(\Dompl(X)\)の
        充満部分圏を\(\mathsf{D}^{+}_{\rr^+}(X)\)で表す．
    \end{enumerate}
\end{DFN}
\end{leftbar}

錐状層がどのように特徴づけられるかを調べる．
次の連続写像を考える．
\[\begin{tikzcd}
    X\ar[r,hook,"j"]&
    X\times\rr^+
    \ar[r,yshift=-0.7ex,"p"']
    \ar[r,yshift=0.7ex,"\mu"]
    &X.
\end{tikzcd}\]
\(j\colon X\to X\times \rr^+\)は
\(x\in X\)を\(X\times\rr^+\)に\((x,1)\)として
埋め込む写像であり．
\(p\colon X\times\rr^+\to X\)は\(X\)への第1射影である．
これらの連続写像を用いて，\(F\in\Dompl(X)\)に対し，次の2つの射を構成する．
\[\begin{tikzcd}
    \mu^{-1}F&
    p^{-1}\Rder{p}_{\ast}\mu^{-1}F
    \ar[l,"\alpha"']
    \ar[r,"\beta"]
    &p^{-1}F.
\end{tikzcd}\]
\(\alpha\)は随伴\((p^{-1},\Rder{p}_{\ast})\)の
余単位\(p^{-1}\Rder{p}_{\ast}\to1_{\Dompl(X)}\)を\(\mu^{-1}F\)に
適用することで得られる．
\(\beta\)は次のように構成される．
\begin{align*}
    p^{-1}\Rder{p}_{\ast}\mu^{-1}F
    &\to
    p^{-1}\Rder{p}_{\ast}\Rder{j}_{\ast}j^{-1}\mu^{-1}F\\
    &\cong
    p^{-1}\Rder(p\circ{j})_{\ast}(\mu\circ{j})^{-1}F\\
    &\cong
    p^{-1}\Rder{1_X}_{\ast}1_X^{-1}F\\    
    &\cong
    p^{-1}F.
\end{align*}

\begin{leftbar}
\begin{PRP}[{\cite[Proposition 3.7.2]{KS90}}]
    \(F\in\mathsf{D}^{+}(X)\)に対し
    次の条件(i)--(v)は同値である．
    \begin{enumerate}[(i)]
        \item \(F\in \mathsf{D}^{+}_{\rr^+}(X)\)
        \item \(\alpha\)と\(\beta\)はどちらも同型である．
        \item すべての\(j\in\zz\)に対し，\(H^j(\mu^{-1}F)\)が\(p\)の各ファイバーで局所定数層となる．
        \item \(p^{-1}F\cong \mu^{-1}F\).
        \item \(p^{!}F\cong \mu^{!}F\).
    \end{enumerate}
\end{PRP}    
\end{leftbar}

\begin{proof}
    次の順に証明する．
    \[\begin{tikzcd}
        (\mathrm{i})\ar[r,Leftrightarrow,"\ref{372-13}"]
        &(\mathrm{iii})\ar[rd,Rightarrow,"\ref{372-32}"']
        &&(\mathrm{iv})\ar[ll,Rightarrow,"\ref{372-43}"']\ar[r,Leftrightarrow,"\ref{372-45}"]
        &(\mathrm{v})\\
        &&(\mathrm{ii})\ar[ru,Rightarrow,"\ref{372-24}"']
    \end{tikzcd}.\]
    
    \begin{enumerate}
        \item\label{372-13} (i) \(\Leftrightarrow\) (iii): 
        まず，\(x\in X\)の\(\rr^+\)軌道は\(
            \mu\left(p^{-1}(x)\right)
        \)と表せることに注意する．実際，\(x\)の\(\rr^+\)軌道\(b\)は
        \[
            b=\left\{\mu(x,t); t\in\rr^+\right\}
        \]であり，これは\(x\)での\(p\)のファイバー\[
            p^{-1}(x)=\left\{(x,t);t\in\rr^+\right\}
        \]の\(\mu\)による像
        \[
            \mu\left(p^{-1}(x)\right)=
            \left\{\mu(x,t);(x,t)\in p^{-1}(x)\right\}
        \]である．
        \begin{align*}
            j_{p^{-1}(x)}&\colon 
            p^{-1}(x)\hookrightarrow X\times \rr^+\\
            j_{\mu\left(p^{-1}(x)\right)}&\colon 
            \mu\left(p^{-1}(x)\right)\hookrightarrow X
        \end{align*}
        をそれぞれ包含写像とすると\[
            j_{\mu\left(p^{-1}(x)\right)}\circ\mu
            =\mu\circ j_{p^{-1}(x)}
        \]が成り立つ．したがって，
        \begin{align*}
            H^j\left(\mu^{-1}F\right)&\cong H^j\left(\mu^{-1}F\right)\rvert_{p^{-1}(x)}\\
            &\cong j_{p^{-1}(x)}^{-1}H^j\left(\mu^{-1}F\right)\\
            &\cong
            j_{p^{-1}(x)}^{-1}\mu^{-1}H^j\left(F\right)\\
            &\cong
            \left(\mu\circ j_{p^{-1}(x)}\right)^{-1}H^j\left(F\right)\\
            &\cong
            \left(
                j_{\mu\left(p^{-1}(x)\right)}\circ\mu
            \right)^{-1}H^j\left(F\right)\\
            &\cong
            \mu^{-1}j_{\mu\left(p^{-1}(x)\right)}^{-1}H^j\left(F\right)\\
            &\cong
            \mu^{-1}\left(H^j\left(F\right)\rvert_{\mu\left(p^{-1}(x)\right)}\right)
        \end{align*}
        である．
        引き戻しが定数層なら元の層も定数層であり，
        定数層の引き戻しも定数層\footnote{
            \(f\colon Y\to X\)を位相空間の間の連続写像とし，
            \(M\)を加群とする．
            \(X\)上の層\(F\)の引き戻し\(f^{-1}F\)が\(Y\)上の
            定数層\(M_Y\)になったとすると，
            \(\mathrm{a}_Y=\mathrm{a}_X\circ f\)より，
            \(
                f^{-1}F
                \cong M_Y
                \cong \mathrm{a}_Y^{-1}M
                \cong f^{-1}\mathrm{a}_X^{-1}M
                \cong f^{-1}M_X
            \)である．
            逆像関手は conservative なので（ホンマか？）\(F\cong M_X\)である．
        }なので，
        \(H^j(\mu^{-1}F)\)が\(p\)の各ファイバーで局所定数層となるのは，
        \(H^j\left(F\right)\)が\(\mu\left(p^{-1}(x)\right)\)で，
        すなわち\(x\)の\(\rr^+\)軌道で局所定数層となるときである．
        \item\label{372-45} (iv) \(\Leftrightarrow\) (v): 
        \(
            p^{-1}F\cong p^{-1}\ltens 
            \omega_{X\times\rr^{+}/X}
        \)と\(
            \mu^{-1}F\cong \mu^{-1}\ltens 
            \omega_{X\times\rr^{+}/X}
        \)が成り立つ．
        \item\label{372-24} (ii) \(\Leftrightarrow\) (iv): 
        \(\alpha\)と\(\beta\)が同型なので，
        互いに逆射となる射\[
            \beta\alpha^{-1}\colon \mu^{-1}F\to p^{-1}F,\quad 
            \alpha\beta^{-1}\colon  p^{-1}F\to \mu^{-1}F
        \]が得られる．したがって，\(p^{-1}F\cong \mu^{-1}F\)である．
        \item\label{372-43} (iv) \(\Leftrightarrow\) (iii): 
        \(\mu^{-1}F\cong p^{-1}F\)とする．
        \(x\in X\)とする．\(
            p^{-1}(x)=\{x\}\times \rr^{+}
        \)に対し\[
            i_{p^{-1}(x)}\colon p^{-1}(x)\hookrightarrow X\times \rr^+,
            \quad
            i_{x}\colon\{x\} \hookrightarrow X
        \]
        とおくと次の図式が可換になる．
        \[\begin{tikzcd}
            p^{-1}(x)
            \ar[r,hookrightarrow, "i_{p^{-1}(x)}"]
            \ar[d, "p\rvert_{p^{-1}(x)}"']
            &X\times \rr^+
            \ar[d, "p"]\\
            \{x\}
            \ar[r,hookrightarrow, "i_{x}"']
            &X
        \end{tikzcd}.\]
        したがって，\(j\)を整数とすると
        \begin{align*}
            H^j\left. \left(p^{-1}F\right)\right\rvert_{p^{-1}(x)}
            &\cong
            p^{-1}\left. H^j\left(F\right)\right\rvert_{p^{-1}(x)}\\
            &\cong
            i^{-1}_{p^{-1}(x)}p^{-1}H^j\left(F\right)\\
            &\cong
            \left(p\circ i_{p^{-1}(x)}\right)^{-1}H^j\left(F\right)\\
            &\cong
            \left(i_{x}\circ p\rvert_{p^{-1}(x)}\right)^{-1}H^j\left(F\right)\\
            &\cong
            \left(p\rvert_{p^{-1}(x)}\right)^{-1}i_{x}^{-1}H^j\left(F\right)\\
            &\cong
            \left(p\rvert_{p^{-1}(x)}\right)^{-1}H^j\left(F\right)_{x}\\
            &\cong
            \left(H^j\left(F\right)_{x}\right)_{p^{-1}(x)}
        \end{align*}
        で，定数層となる（よって特に局所定数層となる）．
        \item\label{372-32} (iii) \(\Leftrightarrow\) (ii):      
    \end{enumerate}
\end{proof}

\begin{leftbar}
\begin{CRL}[{\cite[Corollary 3.7.3]{KS90}}]\label{crl:3.7.3}

\end{CRL}
\end{leftbar}

\begin{leftbar}
\begin{PRP}[{\cite[Proposition 3.7.4]{KS90}}]\label{prop:3.7.4}
    
\end{PRP}
\end{leftbar}

\begin{leftbar}
\begin{PRP}[{\cite[Proposition 3.7.5]{KS90}}]\label{prop:3.7.5}
    \(F\in\Dlcon(E)\)とする．このとき\begin{enumerate}[(i)]
        \item \(\Rder{\tau_{\ast}}F\cong i^{-1}F\).
        \item \(\Rder{\tau_{!}}F\cong i^{!}F\).
    \end{enumerate}
\end{PRP}
\end{leftbar}

\begin{leftbar}
\begin{LMM}[{\cite[Lemma 3.7.6]{KS90}}]\label{lmm:3.7.6}
    
\end{LMM}
\end{leftbar}

\begin{leftbar}
\begin{THM}[{\cite[Theorem 3.7.7]{KS90}}]
    \(\Dlcon(E)\)から\(\Dlcon(E^\ast)\)への
    関手\(\widetilde{\varPhi}_{P'}\)と
    \(\widetilde{\varPsi}_{P}\)は自然に同型である．
\end{THM}
\end{leftbar}

\begin{proof}
    \(
        \widetilde{\varPhi}_{P'}(F)
        =\Rder{p_2}_!\left(p_1^{-1}F\right)_{P'}
    \), \(
        \widetilde{\varPsi}_{P}(F)
        =\Rder{p_2}_{\ast}\RG_{P}(p_1^{-1}F)
    \)なので，これらを同型で結ぶ．
    \begin{align*}
        \widetilde{\varPhi}_{P'}(F)
        &=
        \Rder{p_2}_!\left(p_1^{-1}F\right)_{P'}\\
        &\cong
        \Rder{p_2}_!(\RG_{P}(\left(p_1^{-1}F\right)_{P'}))
        \tag{\(\natural1\)}\label{377-addP}\\
        &\cong
        \Rder{p_2}_!\left(\left(\RG_{P}\left(p_1^{-1}F\right)\right)_{P'}\right)
        \tag{\(\natural2\)}\label{377-changePandP'}\\
        &\cong
        \Rder{p_2}_{\ast}\left(
            \left(\RG_{P}\left(p_1^{-1}F\right)\right)_{P'}
        \right)\tag{\(\natural3\)}\label{377-proper}\\
        &\cong
        \Rder{p_2}_{\ast}\RG_{P}(p_1^{-1}F).\tag{\(\natural4\)}\label{377-removeP'}
    \end{align*}
    ただし，それぞれの同型は以下の通りに示される．
    \begin{description}
        \item[{\eqref{377-addP}}] 完全三角\[
            \Rder{p_2}_!\RG_{P}\left(\left(p_1^{-1}F\right)_{P'}\right)
            \to
            \Rder{p_2}_!\left(p_1^{-1}F\right)_{P'}
            \to
            \Rder{p_2}_!\RG_{E\dpb[Z]E^{\ast}-P}\left(\left(p_1^{-1}F\right)_{P'}\right)
            \to +1
        \]
        において，\(
            \Rder{p_2}_!\RG_{E\tpb[Z]E^{\ast}-P}\left(
                \left(p_1^{-1}F\right)_{P'}
            \right)=0
        \)であることを示せばよい．
        いま命題\ref{prop:3.7.5}より，
        \(i_2\colon E^{\ast}\hookrightarrow E\tpb[Z]E^{\ast}\)を閉うめこみとすると，
        \[
            \Rder{p_2}_!\RG_{E\dpb[Z]E^{\ast}-P}\left(
                \left(p_1^{-1}F\right)_{P'}
            \right)
            \cong
            i_{2}^{!}\RG_{E\dpb[Z]E^{\ast}-P}\left(
                \left(p_1^{-1}F\right)_{P'}
            \right)
        \]
        なので，\begin{align*}
            i_{2}^{!}\RG_{E\dpb[Z]E^{\ast}-P}\left(
                \left(p_1^{-1}F\right)_{P'}
            \right)
            &\cong
            i_{2}^{-1}\RG_{Z\dpb[Z]E^{\ast}}\left(
                \RG_{E\dpb[Z]E^{\ast}-P}\left(
                    \left(p_1^{-1}F\right)_{P'}
                \right)
            \right)\\
            &\cong
            i_{2}^{-1}\RG_{\left(
                Z\dpb[Z]E^{\ast}\cap (E\dpb[Z]E^{\ast}-P)
            \right)}\left(
                \left(p_1^{-1}F\right)_{P'}
            \right)
        \end{align*}
        である．ここで
        \[
            Z\dpb[Z]E^{\ast}\cap (E\dpb[Z]E^{\ast}-P)
            =\left\{(z,0,y);\inner<0,y>=0<0\right\}
        \]
        は空集合であるから，\(\Rder{p_2}_!\RG_{E\tpb[Z]E^{\ast}-P}\left(
                \left(p_1^{-1}F\right)_{P'}
            \right)=0\)である．
        \item[{\eqref{377-changePandP'}}] 演習から従う．
        \item[{\eqref{377-proper}}] 補題\ref{lmm:3.7.6}から\(p_2\)は\(\supp(\RG_{P}\left(p_1^{-1}F\right)_{P'})\)上で固有である．
        \item[{\eqref{377-removeP'}}] 完全三角\[
            \Rder{p_2}_{\ast}\left(
                \left(\RG_{P}\left(p_1^{-1}F\right)\right)_{E\dpb[Z]E^{\ast}-P'}
            \right)
            \to
            \Rder{p_2}_{\ast}\RG_{P}(p_1^{-1}F)
            \to
            \Rder{p_2}_{\ast}\left(
                \left(\RG_{P}\left(p_1^{-1}F\right)\right)_{P'}
            \right)
            \to+1
        \]において，\(
            \Rder{p_2}_{\ast}\left(
                \left(\RG_{P}\left(p_1^{-1}F\right)\right)_{E\dpb[Z]E^{\ast}-P'}
            \right)=0
        \)であることを示せばよい．命題\ref{prop:3.7.5}から
        \[
            \Rder{p_2}_{\ast}\left(
                \left(\RG_{P}\left(p_1^{-1}F\right)\right)_{E\dpb[Z]E^{\ast}-P'}
            \right)
            \cong
            i_{2}^{-1}\left(
                \left(\RG_{P}\left(p_1^{-1}F\right)\right)_{E\dpb[Z]E^{\ast}-P'}
            \right)
        \]
        である．\(\left(\RG_{P}\left(p_1^{-1}F\right)\right)_{E\dpb[Z]E^{\ast}-P'}\)の特徴づけから
        \[\begin{cases}
            \mres[\left(\RG_{P}\left(p_1^{-1}F\right)\right)_{E\dpb[Z]E^{\ast}-P'}]{E\dpb[Z]E^{\ast}-P'}=\mres[\left(\RG_{P}\left(p_1^{-1}F\right)\right)]{E\dpb[Z]E^{\ast}-P'}\\
            \mres[\left(\RG_{P}\left(p_1^{-1}F\right)\right)_{E\dpb[Z]E^{\ast}-P'}]{P'}=0
        \end{cases}\]
        である．これと\(E^{\ast}\cong Z\tpb[Z]E^{\ast}\subset P'\)より，
        \begin{align*}
            i_{2}^{-1}\left(
                \left(\RG_{P}\left(p_1^{-1}F\right)\right)_{E\dpb[Z]E^{\ast}-P'}
            \right)
            &=\mres[\left(
                \left(\RG_{P}\left(p_1^{-1}F\right)\right)_{E\dpb[Z]E^{\ast}-P'}
            \right)]{E^{\ast}}\\
            &=\mres[{\mres[\left(
                \left(\RG_{P}\left(p_1^{-1}F\right)\right)_{E\dpb[Z]E^{\ast}-P'}
            \right)]{P'}}]{E^{\ast}}\\
            &=\mres[0]{E^{\ast}}=0
        \end{align*}
        である．
    \end{description}
\end{proof}
\begin{DFN}[{\cite[Definition 3.7.8]{KS90}}]
    \(F\in\Dlcon(E)\)とする．
    \(F\)の\textbf{フーリエ・佐藤変換} (Fourier-Sato transform) 
    \(F^{\wedge}\)を
    \begin{align*}
        F^{\wedge}
        \coloneqq \widetilde{\varPhi}_{P'}(F)\
        \left(=\Rder{p_2}_!\left(p_1^{-1}F\right)_{P'}\right)\\
        \left(
            \cong\widetilde{\varPsi}_{P}(F)
            =\Rder{p_2}_{\ast}\RG_{P}(p_1^{-1}F)
        \right)
    \end{align*}
    で定める．

    \(G\in\Dlcon(E^\ast)\)とする．
    \(G\)の\textbf{逆フーリエ・佐藤変換} (inverse 
    Fourier-Sato transform) \(G^{\vee}\)を
    \begin{align*}
        G^{\vee}
        \coloneqq \widetilde{\varPsi}_{P'}(F)\
        \left(=\Rder{p_1}_{\ast}\RG_{P'}(p_2^{!}G)\right)\\
        \left(
            \cong \widetilde{\varPhi}_{P}(G)
            =\Rder{p_1}_!\left(p_2^{!}G\right)_{P}
        \right)
    \end{align*}
    で定める．
\end{DFN}

\begin{THM}[{\cite[Theorem 3.7.9]{KS90}}]
    \(\Dlcon(E)\)から\(\Dlcon(E^\ast)\)への
    関手\({}^{\wedge}\)と
    \(\Dlcon(E^\ast)\)から\(\Dlcon(E)\)への
    関手\({}^{\vee}\)は圏同値であり，互いに準逆である．
    とくに，\(F\)と\(F'\)を\(\Dlcon(E)\)の対象とするとき，
    \[
        \Hom_{\Dlcon(E)}(F',F)\cong
        \Hom_{\Dlcon(E)}(F'^{\wedge},F^{\wedge})
    \]
    が成りたつ．
\end{THM}
\begin{LMM}[{\cite[Lemma 3.7.10]{KS90}}]
    \begin{enumerate}[(i)]
        \item \(\gamma\)を\(E\)の固有閉凸錐で零切断を含むものとする．
        このとき，次が成り立つ．\[
            \left(A_\gamma\right)^\wedge
            \cong A_{\Int{\gamma^\circ}}.
        \]
        \item \(U\)を\(E\)の開凸錐とする．このとき次が成り立つ．
        \[
            \left(A_U\right)^\wedge\cong
            A_{{U^{\circ}}^a}\tens \ori_{E^\ast/Z}[-n].
        \]
    \end{enumerate}
\end{LMM}
\begin{RMK}[{\cite[Remark 3.7.11]{KS90}}]
    RMK
\end{RMK}
\begin{PRP}[{\cite[Proposition 3.7.12]{KS90}}]
    \(F\in\Dlcon(E)\)とする．
    \begin{enumerate}[(i)]
        \item \(F^{\wedge\wedge}\cong F^a\tens \ori_{E/Z}[-n]\).
        \item \(U\)を\(E^\ast\)の開凸集合とすると\[
            \RG(U;F^\wedge)
            \cong\RG_{U^\circ}(\tau^{-1}\pi(U);F)
            \cong\RG_{U^\circ}(E;F).
        \]
        \item \(\gamma\)を\(E^\ast\)の固有閉凸錐で
        零切断を含むものとすると\[
            \RG_{\gamma}(E^\ast;F^\wedge)
            \cong
            \RG(\Int{\gamma^{\circ a}};F)\tens\ori_{E/Z}[-n].
        \]
        \item 次が成り立つ．\[
            \left(\rmD'F\right)^\vee\cong\rmD'(F^\wedge),
            \quad
            \left(\rmD F\right)^\vee\cong\rmD(F^\wedge).
            \]
    \end{enumerate}
\end{PRP}
\begin{PRP}[{\cite[Proposition 3.7.13]{KS90}}]
    \begin{enumerate}[(i)]
        \item \(F\in\Dlcon(E)\)とする．このとき次が成り立つ．\[
            (f_\tau^!F)^\wedge\cong f_\pi^!(F^\wedge),
            \quad
            (f_\tau^{-1}F)^\wedge\cong f_\pi^{-1}(F^\wedge).
            \]
        \item \(G\in\Dlcon(E')\)とする．このとき次が成り立つ．\[
            (\Rder{f_\tau}_\ast G)^\wedge\cong \Rder{f_\pi}_\ast(G^\wedge),
            \quad
            (\Rder{f_\tau}_! G)^\wedge\cong \Rder{f_\pi}_!(G^\wedge),
            \]
    \end{enumerate}
\end{PRP}
\begin{PRP}[{\cite[Proposition 3.7.14]{KS90}}]
    \begin{enumerate}[(i)]
        \item \(F\in\Dlcon(E_1)\)とする．このとき次が成り立つ．
        \begin{align*}
            {}^tf^{-1}(F^\wedge)&\cong (\Rder{f}_!F)^\wedge,\\
            {}^tf^{!}(F^\vee)&\cong (\Rder{f}_\ast F)^\vee,\\
            {}^tf^{!}(F^\wedge)&\cong (\Rder{f}_{\ast}F)^\wedge\tens\omega_{E_2^\ast/E_1^\ast},\\
            {}^tf^{!}(F^\vee)&\cong (\Rder{f}_{!}F)^\vee\tens\omega_{E_2^\ast/E_1^\ast}.
        \end{align*}
        \item \(G\in\Dlcon(E_2)\)とする．このとき次が成り立つ．
        \begin{align*}
            (f^{-1}F)^\vee&\cong \Rder^t{f}_!(G^\vee),\\
            (f^{!}F)^\wedge&\cong \Rder^t{f}_\ast(G^\wedge),\\
            (\omega_{E_1/E_2}\tens f^!G)^\vee&\cong \Rder^t{f}_\ast(G^\vee),\\
            (\omega_{E_1/E_2}\tens f^{-1}G)^\wedge&\cong \Rder^t{f}_!(G^\wedge).
        \end{align*}
    \end{enumerate}
\end{PRP}
\begin{PRP}[{\cite[Proposition 3.7.15]{KS90}}]
    \(F_i\in\Dlcon(E_i)\), \(i=1,2\)とする．このとき次が成り立つ．
    \[
        F_1^\wedge\letens[Z]F_2^\wedge
        \cong
        \left(F_1\letens[Z]F_2\right)^\wedge.
    \]
\end{PRP}

\section{その他やったこと}
\begin{proof}[\textbf{\eqref{377-proper}の部分の証明．}]
    \(p_2\)が\(Z\pb[Z]E^{\ast}\)上で固有であることを示せば，
    \(\supp(\RG_{P}(p_1^{-1}F)_{P'})\)で固有となるので，
    \(\Rder{p_2}_!\cong\Rder{p_2}_{\ast}\)となる．
    \(K\subset E^{\ast}\)をコンパクト集合とする．
    \((\mres[p_2]{Z\pb[Z]E^{\ast}})^{-1}(K)\)が
    コンパクトとなることを示す．
    \begin{align*}
        (\mres[p_2]{Z\pb[Z]E^{\ast}})^{-1}(K)
        &=\left\{
            (z,y)\in Z\pb[Z]E^{\ast};
            p_2(z,y)\in K
        \right\}\\
        &=\left\{
            (z,y)\in Z\times K;\id_Z(z)=\pi(y)
        \right\}\\
        &=\left\{
            (z,y)\in Z\times K;z=\pi(y)
        \right\}\\
        &=\pi(K)\times K
    \end{align*}
    である．
    連続写像によるコンパクト集合の像は
    コンパクトなので\(\pi(K)\)はコンパクトである．
    さらに，コンパクト集合の積はコンパクトなので\(\pi(K)\times K\)は
    コンパクトである．
\end{proof}

\begin{PRP}
    \(X\)を局所コンパクト空間とし，\(F\)を\(X\)上の層とする．
    次の同型が成り立つ．
    \[
        \indlim[K\subset X]H_K^j(X;F)
        \cong
        H^j_c(X;F).
    \]
\end{PRP}
\begin{proof}
    \(F\)の脆弱層による分解\(F\qis I^\bullet\)を取れば，
    \(\RG_K(X;F)^i\cong\Gamma_K(X;I^i)\), 
    \(\RG_c(X;F)^i\cong\Gamma_c(X;I^i)\)である．
    したがって，一般に層\(F\in\Sh(X)\)に対して\[
        \indlim[K\subset X]\Gamma_K(X;F)\cong \Gamma_c(X;F)
    \]が成り立つことを示せばよい．
\end{proof}










%\section{aa}
%===============================================
% 参考文献スペース
%===============================================
\begin{thebibliography}{20} 
  \bibitem[dR]{dR} 
  de Rham, \textit{Vari\'et\'es diff\'erentiables}, Hermann, 1953.
  和訳：ド・ラーム, 微分多様体, 東京図書, 1974.
  \bibitem[GKS12]{GKS12} St\'ephane Guillermou, 
  Masaki Kashiwara, Pierre Schapira, 
  \textit{Sheaf Quantization of Hamiltonian Isotopies and Applications to Nondisplaceability Problems}, 
  Duke. Math. J. 161, No. 2, pp.201--245, 2012.  
  \bibitem[GP74]{GP74} Victor Guillemin, Alan Pollack, 
    \textit{Differential Topology}, 
    Prentice-Hall, 1974.
  \bibitem[KS90]{KS90} Masaki Kashiwara, Pierre Schapira, 
    \textit{Sheaves on Manifolds}, 
    Grundlehren der Mathematischen Wissenschaften, 292, Springer, 1990.
  \bibitem[Hat89]{Hat89} 服部晶夫, 多様体 増補版, 岩波全書 288, 岩波書店, 1989.
  %\bibitem[Og02]{Og02} 小木曽啓示, 代数曲線論, 朝倉書店, 2022.
  \bibitem[Hor63]{Hor63} Lars H\"ormander, 
  \textit{Linear Pertial Differential Operators}, Fourth Printing, 
  Grundlehren der Mathematischen Wissenschaften, 116, Springer, 1963.
  \bibitem[Hor71]{Hor71} Lars H\"ormander, 
  \textit{Fourier Integral Operators I}, 
  Acta Math. 127, 79--183, (1971).
  \bibitem[Hor89]{Hor89} Lars H\"ormander, 
  \textit{The Analysis of Linear Pertial Differential Operators I}, Second Edition,
  Grundlehren der Mathematischen Wissenschaften, 256, Springer, 1989.
  \bibitem[Hor85]{Hor85} Lars H\"ormander, 
  \textit{The Analysis of Linear Pertial Differential Operators III}, 
  Grundlehren der Mathematischen Wissenschaften, 274, Springer, 1985.
  \bibitem[Roc]{Roc} Rockaffeller, \textit{Convex Analysis},. 
  \bibitem[Sch85]{Sch85} Pierre Schapira, 
    \textit{Microdifferential Systems in the Complex Domain}, 
    Grundlehren der Mathematischen Wissenschaften, 269, 
    Springer, 1985.
    \bibitem[ST92]{ST92} Pierre Schapira, Nobuyuki Tose, 
    \textit{Morse Inequalities for \(\rr\)-constructible Sheaves}, 
    Adv. Math. 93, pp.1--8, 1992. 
    \bibitem[Zwo12]{Zwo12} Maciej Zworski,
    \textit{Semiclsaaical Analysis}, 
    Graduate Studies in Mathematics, 138, 
    American Mathematical Society, 
    2012. 
\end{thebibliography}

%===============================================




\end{document}
